% !TeX program = xelatex
\documentclass[UTF8,a4paper,zihao=-4,fontset=windows,oneside,openany]{ctexbook}
\usepackage[margin=22mm,headheight=15pt]{geometry}
\usepackage{amsmath,amssymb,graphicx,xcolor,needspace,longtable,fancyhdr}
\usepackage[unicode,hidelinks]{hyperref}
\setlength{\parindent}{0pt}
\setlength{\parskip}{4pt}
\setlength{\emergencystretch}{3em}
\linespread{1.2}
\raggedbottom
\pagestyle{fancy}
\fancyhf{}
\fancyhead[L]{湘教版数学必修第一册：教师转写工作稿}
\fancyhead[R]{题干与公式可编辑}
\fancyfoot[C]{\thepage}
\renewcommand{\headrulewidth}{0.3pt}
\newcommand{\subq}[1]{\par\Needspace{2\baselineskip}\noindent（#1）\quad}
\newcommand{\choice}[1]{\par\noindent #1.\quad}
\newcommand{\blank}{\underline{\hspace{1.7cm}}}
\newcommand{\ocrglyph}[1]{\textcolor{red!70!black}{[识别字符 U+#1 待核]}}
\newenvironment{bankitem}[4]{\par\Needspace{5\baselineskip}\noindent
  {\bfseries #1}\par{\small\color{blue!55!black}#2}\par
  {\footnotesize\ttfamily\detokenize{#3}}\par
  {\footnotesize\color{red!65!black}#4}\par\nopagebreak}
  {\par\medskip\hrule\medskip}
\newenvironment{ocrdraft}{\begingroup\small\color{black!65}}{\par\endgroup}
\newcommand{\questioninput}[1]{\input{questions/#1.tex}}
\hypersetup{pdftitle={湘教版必修第一册教材题库：转写工作稿},pdfauthor={教学资料整理}}
